We have presented DialogGSR, a dialog generation model with generative subgraph retrieval. 
DialogGSR retrieves context-relevant subgraphs, by generating the subgraph token sequences considering both the dialog context and the graph information.
We have proposed novel knowledge graph linearization to convert knowledge triplets into token sequences with self-supervised graph-specific tokens to represent knowledge graphs without separate knowledge graph modules.
In addition, we have formulated a graph-constrained decoding for valid and relevant generative retrieval.
Our experiments demonstrate the effectiveness of our proposed method in knowledge\textendash{graph} grounded dialog generation.
Our codes are publicly available at \url{https://github.com/mlvlab/DialogGSR}.
\section*{Limitations}
The proposed DialogGSR generatively retrieves token sequences of the subgraph from a knowledge graph and then generates a response with the retrieved subgraph.
However, similar to works using graph retrieval on knowledge-grounded dialog generation, our generative subgraph retrieval can retrieve only the knowledge information contained in the knowledge graph.
Second, the benchmark datasets for knowledge graph{\textendash}grounded dialog generation are limited.
Therefore, new benchmark datasets on dialog generation with knowledge graphs warrants greater attention.


\section*{Ethics Statement}
Our DialogGSR does not have any direct negative social impacts, but it can potentially be used maliciously, similar to other dialog generation models. These models may produce factually incorrect or biased responses, particularly in sensitive areas such as politics, religion, and diplomacy. To address these risks, we advocate for the release of benchmark datasets without private information and emphasize the need for research into the methods that detect the source of texts. These measures are essential for the responsible development and use of dialog generation technologies.
